\section{Risk models: factor exposures and their control}

\begin{qfmodel}[Factor risk model: factor returns, exposures, residuals and the factor covariance, with economic and statistical factors and their use in limiting unintended exposures]\label{mod:factor-risk-model}
\qffield{Statement} In Narang's framing a \emph{risk} is an exposure that accompanies the pursuit of alpha but is not intentionally forecast: ``risk exposures are those that are not intentionally sought out by the nature of whatever forecast the quant is making in the alpha model''. A risk model measures such exposures, highlights them and provides input to the portfolio construction model, which decides whether to accept, limit or eliminate each \citep[chap.~4, ``Risk Models'', pp.~71--73]{s7-narang}. Risk \emph{management}, in his words, ``is about the intentional selection and sizing of exposures to improve the quality and consistency of returns'' (p.~71). The factor model that underlies most risk models, as Boyd, Johansson, Kahn, Schiele and Schmelzer write it, decomposes the $n$ asset returns into $k$ factor returns and idiosyncratic returns \citep[author version of 5 Jan.\ 2024, sec.~3.3, eq.~(14), pp.~20--21]{wide-boyd-markowitz-seventy}:
\[ r=Ff+\epsilon,\qquad \Sigma=F\Sigma_fF^\top+D,\qquad \sigma^2=(F^\top w)^\top\Sigma_f(F^\top w)+w^\top Dw, \]
with $F$ the $n\times k$ loading (exposure) matrix, known when the portfolio is built, $f$ the factor returns with covariance $\Sigma_f$, and $\epsilon$ the idiosyncratic returns with diagonal covariance $D$ (positive entries), uncorrelated with $f$. The means of $\epsilon$ are the assets' alphas.

\qffield{Notation} $F$, $f$, $\Sigma_f$, $\epsilon$, $D$, $w$ as in Boyd et al.; $F^\top w$ is the portfolio's vector of factor exposures; $\sigma^2$ the portfolio variance. Usually $k\ll n$, so $\Sigma$ is low rank plus diagonal.

\qffield{Conditions} \emph{Assumptions.} The decomposition is exact only under its assumptions: factor and idiosyncratic returns uncorrelated, and idiosyncratic returns uncorrelated across assets ($D$ diagonal). Interpretation: under them, an equally weighted portfolio of $n$ assets with unit loading on one factor of 15 percent volatility and residuals of 30 percent has volatility 33.5, 17.8, 15.3 and 15.0 percent for $n=1$, 10, 100 and 1{,}000, so the residual part diversifies away and the factor part does not. This is what Narang means by systematic risks that ``cannot be diversified away'' (p.~77). A diagonal $D$, with bounded residual variances and equal weights, is sufficient for this, not necessary. What matters is how much residual covariance remains when it is summed over pairs (Interpretation).

With a constant residual correlation of 0.2 between every pair and residual variance 0.09, equal-weighted residual variance is $0.09\,(0.2+0.8/n)=0.018+0.072/n$, which stays above 0.018 however large $n$ is: that common correlation behaves like an omitted factor. With the same correlation of 0.2 only within disjoint pairs of assets, it is $0.108/n$, which still vanishes as $n$ grows. In these examples correlation confined to small groups still diversifies, while a correlation shared by every pair leaves a floor.

\emph{Two kinds of factor.} Theory-driven or economic factors are named in advance on an economic argument: for equities the market, sectors and market capitalisation; for bonds, interest rates (pp.~77--78). Boyd et al.\ give market, industry and style portfolios such as Fama--French for equities and yield-curve, rate and spread factors for bonds, which are ``interpretable''. Empirical or statistical factors come from historical returns by methods such as principal component analysis; apart from the first, ``which typically is close to the market return, these factors are less interpretable'' (Boyd et al., p.~20). Narang says the first principal component of Treasury yields usually corresponds to the level of rates, and in equities the market comes first and then sectors (p.~78).

Interpretation: a principal component is defined only up to sign and can change order and composition from one sample to the next, so its economic label is an after-the-fact reading, not a property of the method. Narang notes that statistical models can be ``fooled'' into factors that do not persist or into spurious exposures (p.~78).

\qffield{Practice} \emph{Uses} (pp.~72 and~82--83): measuring which systematic exposures a portfolio takes and how large they are; explaining profits and losses; and, inside an alpha model, forecasting the residuals of returns against the risk model's factors. \emph{Limiting exposures.} Narang's general rule is ``to eliminate any unintentional exposures, since there should be no expectation of being compensated sufficiently for accepting them'' (p.~76); his examples are a value strategy's incidental sector and small-cap bets (pp.~72 and~77) and a view that CVX would beat XOM, expressed as a long-only CVX trade that also bets on the market, which buying CVX and shorting an equivalent amount of XOM ``substantially eliminated'' (p.~76).

Interpretation: equal dollar amounts hedge the market exactly only if the two stocks have equal market betas; otherwise the pair keeps a net market exposure of the beta difference times the dollar size, so the hedge is approximate, and exact neutrality needs beta-weighted sizes under a stated risk model. Exposure limits, such as a $\pm5$ percent net market exposure, apply to risks the alpha model does not forecast (p.~75). He prefers, ``to the extent possible'', to build exposure decisions into the alpha model's bet structure, for example forecasting stocks relative to their sectors, rather than relying on one risk model in an optimiser (pp.~76--77 and~80--81).

Boyd et al.\ express neutrality to factor $i$ as zero covariance of the portfolio with that factor, $(\Sigma_fF^\top w)_i=0$, a linear constraint (sec.~3.3, eq.~(15), p.~21). Interpretation: with correlated factors this differs from zero exposure, $(F^\top w)_i=0$. For three assets with loadings $(1,0.5)$, $(1,-0.5)$, $(0.8,1)$, factor variances 0.04 and 0.01 and factor covariance 0.012, the weights $(0.5,-0.5,0)$ have exposures $(0,0.5)$ but a covariance of 0.006 with the first factor.

\emph{Choosing a model.} Most quants use theory-driven models for their interpretability; empirical ones promise adaptiveness but lag a regime change and trade adaptiveness against estimation error. With two years of daily data each new day changes two of about 520 observations, while shorter windows lose statistical power. Intraday data is one partial remedy. The two kinds can be combined, and commercial models (BARRA, Northfield, Qontigo/Axioma, Wolfe) are mostly theory-driven and for equities (pp.~78--80). Boyd et al.\ add that a factor model can estimate a large covariance better than direct estimation and cuts the solve cost from $O(n^3)$ to $O(nk^2)$ (p.~21).

\qffield{Pitfalls} \emph{Risk model against decision.} Narang assigns the intentional selection and sizing of exposures to risk \emph{management}, not to the risk model; in his black-box schematic (Exhibit~4.1, p.~83) the risk model is an input to the portfolio construction model, which selects and sizes. \emph{Factor, exposure, signal or premium?} Narang's risk chapter names market, sector and market-capitalisation risk, not market, size, value and momentum as ``the most common factors''; value is his example of an alpha (pp.~71--72). Whether such a factor is a risk to limit, a signal to forecast or a premium expected to be paid depends on the strategy's forecasts, not on the factor.

The same factor can play several roles. \emph{Unintended is not the same as uncompensated.} Narang's rule is that ``there should be no expectation of being compensated sufficiently'' for unintentional exposures (p.~76): not that they earn nothing, but that any expected reward is not enough for the risk. Whether a given exposure is compensated, and by how much, is an empirical question about risk premia, not settled by its being unintended (Interpretation).

\emph{Empirical, not adaptive by default.} Updating an empirical model with new data is not evidence that its forecasts improve; Narang himself describes the lag and the loss of power (p.~79). His remark that empirical models are ``most commonly found among'' statistical arbitrage traders (p.~78) describes their users; it is not a definition.

\qffield{See also} \hyperref[mod:mean-variance-portfolio-optimisation]{mean--variance optimisation} (where exposure limits, penalties and the covariance enter); \hyperref[meth:rule-based-position-sizing]{rule-based position sizing}; \hyperref[def:quant-strategy-risk-sources]{sources of strategy risk} (regime change and model risk); \hyperref[def:common-investor-contagion]{common-investor contagion} (VaR-based leverage under rising volatility).

\qffield{Further reading} Covariance shrinkage is in \hyperref[meth:portfolio-covariance-estimation]{covariance estimation}; estimating factor-model loadings requires a separate estimation procedure.

\end{qfmodel}

\begin{qfmethod}[Estimating a large covariance matrix: why the sample matrix fails when assets are many relative to observations, and linear and nonlinear shrinkage]\label{meth:portfolio-covariance-estimation}
\qffield{Statement} Ledoit and Wolf's review starts from the case ``when the number of variables is of the same magnitude as the number of observations'', where ``the sample covariance matrix certainly will not do'' \citep[UZH ECON WP~323, version of February 2020, abstract, p.~1]{wide-ledoit-wolf-review}. \emph{Linear shrinkage} to the identity solves $\min_{\rho_1,\rho_2}E\|\hat\Sigma-\Sigma\|_F^2$ over $\hat\Sigma=\rho_1I+\rho_2S$, with the scaled Frobenius norm $\|A\|_F^2=\operatorname{Tr}(A'A)/N$ (eqs.~(2.1)--(2.5), pp.~4--5). The solution is
\[ \Sigma^*=\gamma^*\mu I+(1-\gamma^*)S,\qquad \gamma^*=\frac{\beta^2}{\delta^2}, \]
with $\mu=\langle\Sigma,I\rangle$ the average variance, $\alpha^2=\|\Sigma-\mu I\|_F^2$, $\beta^2=E\|S-\Sigma\|_F^2$ and $\delta^2=E\|S-\mu I\|_F^2=\alpha^2+\beta^2$ (eqs.~(2.6)--(2.7), p.~5). It is a convex combination, positive definite and invertible ``even when $N>T$''. The optimal intensity depends on unknown population quantities, so a feasible estimator is derived under large-dimensional asymptotics, with $N/T\to c$ (sec.~2.2, p.~6).
\begin{itemize}
\item \emph{Other targets.} Linear shrinkage can also target the single-factor matrix or the constant-correlation matrix.
\item \emph{Their own comparison.} In their 2004 back-test of a tracking-error-constrained manager, the single-factor target, the constant-correlation target and a five-principal-component factor model all ``dominated the sample covariance matrix, but there was no clear winner'': constant correlation was best for $N\le100$ and the single factor for $N\ge225$ (sec.~3.3, p.~14).
\item \emph{Nonlinear shrinkage} keeps the sample eigenvectors and shrinks each eigenvalue differently. It ``can deliver another level of performance improvement, especially if overlaid with'' time-varying co-volatility or factor models (abstract; secs.~4--6).
\end{itemize}

\qffield{Notation} $N$ assets, $T$ observations, $S$ sample covariance, $c=N/T$ the concentration; the norm divides by $N$, unlike the Frobenius norm.

\qffield{Conditions} Section~2 assumes i.i.d., mean-zero observations; demeaned data and dynamic (DCC) and factor extensions are treated later (Remark~2.1; secs.~5--6). Optimality is in Frobenius loss, which is not the same as the out-of-sample variance of an optimised portfolio, though the review's applications are mostly Markowitz problems.

\qffield{Practice} Own simulation with a known one-factor population of 100 assets and 120 observations:
\begin{itemize}
\item $\alpha^2+\beta^2=\delta^2$ holds (0.2035 against 0.2031);
\item the oracle intensity is 0.045, and the Frobenius loss falls from 0.00910 to 0.00871;
\item the true variance of the global minimum-variance portfolio falls from 0.0252 using $S$ to 0.0089 using $\Sigma^*$, against an optimum of 0.0044.
\end{itemize}
The small loss gain and the large portfolio gain illustrate that the two criteria differ. Narang's grouping of similar instruments (\hyperref[mod:factor-risk-model]{factor risk model}) is a structured target of this kind, without a rule for the intensity.

\begin{itemize}
\item \emph{The constant-correlation target and its feasible intensity} (\nolinkurl{s6-honey}, 2003 manuscript). The target $F$ keeps the sample variances and sets every covariance to $\bar r\sqrt{s_{ii}s_{jj}}$, with $\bar r$ the average sample correlation (App.~A, eq.~(3), p.~12). Under Frobenius loss the optimal weight on $F$ behaves like $\kappa/T$ with $\kappa=(\pi-\rho)/\gamma$: $\pi$ sums the asymptotic variances of the $\sqrt T s_{ij}$, $\rho$ the asymptotic covariances of $\sqrt Tf_{ij}$ with $\sqrt Ts_{ij}$, and $\gamma=\sum_{ij}(\phi_{ij}-\sigma_{ij})^2$ measures the target's misspecification. Each is replaced by a consistent sample estimate, and the intensity is $\hat\delta^*=\max\{0,\min\{\hat\kappa/T,1\}\}$ (App.~B, eqs.~(4)--(5), pp.~12--15) \citep[App.~A--B, pp.~12--15]{s6-honey}.

Numerically, with 50 assets and 60 observations the shrunk matrix has the lower Frobenius loss; the intensity goes to one when the population really has a constant correlation (the target is then right) and is about 0.3 when correlations come in blocks; with more assets than observations the shrunk matrix is invertible and the sample matrix is not.
\item \emph{POET} (\nolinkurl{wide-fan-liao-mincheva-poet}, arXiv 1201.0175). Keep the first $K$ principal components of the sample covariance and threshold the rest, the ``principal orthogonal complement'': off-diagonal entries $\hat r_{ij}$ of the remainder are set to $s_{ij}(\hat r_{ij})I(|\hat r_{ij}|\ge\tau_{ij})$ with an adaptive threshold $\tau_{ij}=\tau(\hat r_{ii}\hat r_{jj})^{1/2}$, so $\tau=0$ gives the sample covariance and $\tau=1$ a strict factor model (sec.~2, eqs.~(2.4)--(2.7), printed pp.~9--10). It assumes \emph{pervasive} factors (the eigenvalues of $p^{-1}B'B$ bounded away from zero, so the first $K$ eigenvalues grow with the number of assets $p$) and a sparse residual covariance \citep[sec.~2, printed pp.~6--10]{wide-fan-liao-mincheva-poet}.
\item \emph{Dynamic correlation} (Engle and Sheppard's account of DCC, NBER WP~8554, 2001; Engle's 2002 \emph{JBES} paper is \nolinkurl{wide-engle-dcc}). Each series gets its own univariate GARCH; the standardised residuals $\epsilon_t$ drive $Q_t=(1-\sum\alpha_m-\sum\beta_n)\bar Q+\sum\alpha_m\epsilon_{t-m}\epsilon'_{t-m}+\sum\beta_nQ_{t-n}$, with $\bar Q$ the unconditional covariance of the standardised residuals, and the correlation matrix is $R_t=Q_t^{*-1}Q_tQ_t^{*-1}$ with $Q_t^*$ the square roots of $Q_t$'s diagonal; $\alpha_m,\beta_n\ge0$ with $\sum\alpha_m+\sum\beta_n<1$ and a positive-definite $\bar R$ keep $R_t$ a correlation matrix (eqs.~(1)--(2) and Proposition~2, printed pp.~5--6) \citep[secs.~2--3, printed pp.~5--6]{engle-sheppard-2001-dcc}. Estimation is in two steps, the variances first and then the correlation parameters.
\item \emph{The published review} (\emph{Journal of Financial Econometrics} 20(1), 2022). As a rule of thumb, ``one generally needs both $N$ and $T$ to be above 50 for nonlinear shrinkage to outperform linear shrinkage'' when the optimal transformation is markedly nonlinear, and then ``there is basically nothing to lose but potentially a lot to gain by upgrading'' (sec.~3.3, pp.~205--206). For Markowitz problems with time-varying co-volatility they propose DCC-NL: univariate GARCH(1,1) to devolatilise, nonlinear shrinkage of the unconditional correlation matrix, and composite likelihood for the dynamics; RiskMetrics 2006 ``performed poorly and cannot be recommended for the purpose of Markowitz's portfolio selection'' (sec.~4, pp.~208--210, n.~13) \citep[2022 journal version, secs.~3.3--4, pp.~205--210]{wide-ledoit-wolf-review}. The working paper read earlier and the journal version agree on the passages compared (abstract and the linear-shrinkage solution).
\end{itemize}

\qffield{Pitfalls} The nonlinear estimator itself was not transcribed. For how the loadings, factor covariance and residual variances of a factor model are estimated, see Connor, Goldberg and Korajczyk (\nolinkurl{wide-connor-goldberg-korajczyk}). \emph{Time-varying estimators}. Boyd et al.\ list EWMA, Engle's DCC and iterated EWMA as approaches that ``work well in practice'', with or without a factor model. They give no comparison and refer to Johansson et al.\ (2023) for detail \citep[arXiv v1, sec.~3.4, pp.~21--22]{wide-boyd-markowitz-seventy}. In their backtest the estimator is $\hat\Sigma_t=\alpha_t\sum_{\tau=1}^t\beta^{t-\tau}r_\tau r_\tau^\top$ with $\alpha_t=(1-\beta)/(1-\beta^t)$. It uses the second moment, the mean being negligible, with a half-life of 125 trading days, $\beta\approx0.994$ ($0.5^{1/125}=0.9945$). The authors note that the half-life does not change their results qualitatively (pp.~31--32). Interpretation: the normalisation makes the weights sum to one in short samples, which a fixed $1-\beta$ does not. The univariate GARCH and EWMA comparison is in \hyperref[mod:garch-volatility-forecast]{GARCH volatility forecasts}.

\qffield{See also} \hyperref[mod:factor-risk-model]{factor risk model}; \hyperref[mod:mean-variance-portfolio-optimisation]{mean--variance optimisation}.

\qffield{Further reading} \nolinkurl{wide-connor-goldberg-korajczyk}.

\end{qfmethod}

\begin{qfevidence}[How much of an empirical correlation matrix is noise: the Marchenko--Pastur band against 406 S\&P~500 stocks]\label{ev:correlation-matrix-noise}
\qffield{Statement} For a $T\times N$ matrix of independent entries of variance $\sigma^2$, with $Q=T/N\ge1$ fixed as both grow, the eigenvalue density of the correlation matrix is $\rho(\lambda)=\frac{Q}{2\pi\sigma^2}\frac{\sqrt{(\lambda_{\max}-\lambda)(\lambda-\lambda_{\min})}}{\lambda}$ on $[\lambda_{\min},\lambda_{\max}]$, with $\lambda_{\max,\min}=\sigma^2(1+1/Q\pm2\sqrt{1/Q})$ \citep[arXiv cond-mat/9810255v1, eq.~(3), p.~1]{wide-laloux-noise-dressing}. Laloux, Cizeau, Bouchaud and Potters compare it with the correlation matrix of $N=406$ S\&P~500 stocks over $T=1309$ days, 1991--96 ($Q=3.22$). The largest eigenvalue, the market mode, is ``25 times larger than the predicted $\lambda_{\max}$''. Taking the rest as noise and fitting $\sigma^2$, the best fit ($\sigma^2=0.74$) ``accounts quite satisfactorily for 94\% of the spectrum, while the 6\% highest eigenvalues still exceed the theoretical upper edge'' (pp.~1--2).

\qffield{Notation} $N$ assets, $T$ observations, $Q=T/N$; returns demeaned and scaled to unit volatility before the correlation matrix is formed.

\qffield{Conditions} Evidence on 1991--1996 data. The band is an $N\to\infty$ limit for independent, identically distributed noise; for finite $N$ the edges blur. Heavy tails and volatility clustering in returns also move the edges, which the authors say a smaller effective $Q$ could absorb (p.~2). Eigenvalues inside the band are not proved to be noise, only indistinguishable from it on this test.

\qffield{Practice} With pure noise at the paper's $N$ and $T$, the simulated eigenvalues fall within the band $[0.196,2.424]$ up to the finite-sample blur. \emph{A slip in the source} (own reading): the text derives the residual variance as ``$\sigma^2=1-\lambda_{\max}/N=0.85$''. With $\lambda_{\max}=2.43$ that would give 0.994; with the market eigenvalue $\lambda_1=25\lambda_{\max}\approx60.6$, $1-\lambda_1/406=0.851$, so the formula must mean $\lambda_1$. The practical use is in \hyperref[meth:portfolio-covariance-estimation]{covariance estimation}: clipping or shrinking the eigenvalues inside the band, which Bun, Bouchaud and Potters review (\nolinkurl{wide-bun-bouchaud-potters-cleaning}).

\qffield{Pitfalls} The sample is 1991--96 daily data for one index; the 94/6 split is a fit with $\sigma^2$ chosen to match, not a test. Locators refer to the arXiv preprint of the letter (\emph{Physical Review Letters} 83, 1999).

\qffield{See also} \hyperref[meth:portfolio-covariance-estimation]{covariance estimation}; \hyperref[mod:pca-risk-distribution-weights]{PCA risk-distribution weights}.

\qffield{Further reading} \nolinkurl{wide-bun-bouchaud-potters-cleaning}; \nolinkurl{wide-vershynin-hdp}.

\end{qfevidence}

\begin{qfmodel}[GARCH(1,1) and EWMA (RiskMetrics) variance forecasts: mean reversion against a flat forecast, and the evidence on forecast accuracy]\label{mod:garch-volatility-forecast}
\qffield{Statement} The GARCH(1,1) daily variance is $\sigma^2_{t+1}=\omega+\alpha R_t^2+\beta\sigma_t^2$ with $\alpha+\beta<1$. Its long-run variance is $\sigma^2=\omega/(1-\alpha-\beta)$, so $\sigma^2_{t+1}=\sigma^2+\alpha(R_t^2-\sigma^2)+\beta(\sigma_t^2-\sigma^2)$, and the $k$-day-ahead forecast reverts at the rate of the \emph{persistence} $\alpha+\beta$ \citep[chap.~4, secs.~2--3, printed pp.~70--71 of an uncorrected proof]{wide-christoffersen-elements}:
\[ E_t[\sigma^2_{t+k}]-\sigma^2=(\alpha+\beta)^{k-1}(\sigma^2_{t+1}-\sigma^2). \]
RiskMetrics' exponential smoothing, $\sigma^2_{t+1}=\lambda\sigma^2_t+(1-\lambda)R_t^2$ with $\lambda=0.94$ set for every asset, is the case $\omega=0$, $\alpha+\beta=1$. It ``ignores the fact that the long-run average variance tends to be relatively stable over time'', so its forecast is flat in $k$ and ``counterfactual'' at longer horizons. Its advantages are a single parameter, no estimation and little stored data (pp.~70--71).

\emph{Evidence on accuracy.} Hansen and Lunde compare 330 volatility models out of sample against intraday realised volatility: the DM--dollar rate (1,254 daily returns, October 1987 to September 1992) and IBM (2,378, January 1990 to May 1999). They use seven loss functions and White's and Hansen's tests for superior predictive ability. They find ``the best models do not provide a significantly better forecast than the GARCH(1,1) model'', while ``if an ARCH(1) model is selected as the benchmark, it is clearly outperformed'' \citep[abstract; sec.~2.4, pp.~7--8; sec.~6]{hansen-lunde-2001}. Their model set includes IGARCH, which they find ``surprisingly poor'' on most criteria (p.~25), but not RiskMetrics' fixed-$\lambda$ smoother as such.

\emph{Two extensions}.
\begin{itemize}
\item \emph{Asymmetry (GJR).} Glosten, Jagannathan and Runkle let the squared innovation enter the variance with a different coefficient according to its sign: $v_{t-1}=b_0+b_1v_{t-2}+g_1\epsilon^2_{t-1}+g_2\epsilon^2_{t-1}I_{t-1}$ (Model~2, eq.~(9), printed p.~10) \citep[Minneapolis Fed Working Paper~505, November 1992, printed p.~10]{wide-gjr-garch}. In this working paper $I_{t-1}=1$ when $\epsilon_{t-1}$ is \emph{positive}; the common textbook form puts the indicator on negative shocks, so the sign of the asymmetry coefficient depends on the convention used.

On monthly U.S.\ excess returns they find that ``positive unanticipated returns result in a downward revision of the conditional volatility whereas negative unanticipated returns result in an upward revision \dots\ of a similar magnitude'', that volatility is less persistent than the standard model suggests, and a negative relation between conditional mean and variance once the model allows for asymmetry, seasonality and interest rates (abstract). The paper was published in the \emph{Journal of Finance} (1993).
\item \emph{Realised-volatility cascade (HAR-RV).} Corsi models next-day realised volatility as $RV^{(d)}_{t+1}=c+\beta^{(d)}RV^{(d)}_t+\beta^{(w)}RV^{(w)}_t+\beta^{(m)}RV^{(m)}_t+\omega_{t+1}$, where $RV^{(w)}$ and $RV^{(m)}$ average the daily values over 5 and 22 trading days (eqs.~(4) and~(6)--(8), pp.~177--181) \citep[\emph{Journal of Financial Econometrics} 7(2), 2009, pp.~174--196]{wide-corsi-har}. It is a restricted AR(22) estimable by least squares, which reproduces long memory approximately without being a long-memory process; Corsi's $RV$ is realised \emph{volatility}, the square root of summed squared intraday returns, and his averages are of volatilities, which by Jensen's inequality are not exactly the volatility of the longer period (n.~4). The HAR forecast equals the AR(22) with coefficients $\beta^{(d)}+\beta^{(w)}/5+\beta^{(m)}/22$ at lag~1, $\beta^{(w)}/5+\beta^{(m)}/22$ at lags 2--5 and $\beta^{(m)}/22$ at lags 6--22.
\end{itemize}

\qffield{Notation} $R_t$ daily return (mean taken as zero); $\sigma^2_{t+1}$ the variance for day $t+1$ known at the end of day $t$; $\lambda$ the smoothing weight.

\qffield{Conditions} The GARCH forecast formula needs $\alpha+\beta<1$ and parameters estimated by quasi-maximum likelihood. Comparisons depend on the volatility proxy: Hansen and Lunde argue that noisy proxies such as squared daily returns obscure the differences between models (sec.~1). They also depend on the loss function and the sample, here one currency and one stock in the 1990s.

\qffield{Practice} Synthetic check. The $k$-step formula matches simulation. Where GARCH(1,1) is true (persistence 0.98), the EWMA forecast's mean squared error for cumulative variance is 0.9 percent above GARCH's at one day and 15 percent above at 20 days. That shows how the flat forecast costs more as the horizon grows; it is not evidence on market data. This bears on Narang's statements (\hyperref[mod:mean-variance-portfolio-optimisation]{mean--variance optimisation}). That GARCH is ``the most widely used such technique (particularly in equity trading)'' is unsupported, since no source here measures usage. That a simpler EWMA ``may perform just about as well'' is consistent with the check at short horizons, but Hansen and Lunde did not test the fixed-$\lambda$ smoother.

\qffield{Pitfalls} The GARCH formulation follows secondary sources, not Bollerslev's 1986 paper (\nolinkurl{s7-bollerslev}). The Christoffersen file is an uncorrected proof of the second edition, so its page numbers may differ from print. The Hansen--Lunde text is their March 2001 working paper (CAF Working Paper Series, No.~84); the published version (\emph{Journal of Applied Econometrics} 20(7), 2005) may summarise the IBM results differently.

\qffield{See also} \hyperref[meth:var-es-estimation-methods]{VaR and ES estimation} (filtered historical simulation scales by a GARCH variance); \hyperref[meth:portfolio-covariance-estimation]{covariance estimation}; \hyperref[meth:multiple-testing-control]{multiple testing} (White's Reality Check, which Hansen and Lunde use).

\qffield{Further reading} \nolinkurl{s7-bollerslev}; \nolinkurl{s7-realized} (realised volatility); \hyperref[meth:volatility-forecast-loss]{loss functions for comparing volatility forecasts}.

\end{qfmodel}

\begin{qfmethod}[Comparing volatility forecasts with a noisy proxy: which loss functions keep the ranking, and why MSE and QLIKE do while MAE does not]\label{meth:volatility-forecast-loss}
\qffield{Statement} True conditional variance is unobserved, so forecasts are compared against a proxy $\hat\sigma^2_t$ (a squared return, a range, realised variance) that is conditionally unbiased but noisy. Patton calls a loss function \emph{robust} if ranking two forecasts by expected loss gives the same order with the proxy as with the true variance (Definition~1, printed p.~4). A loss is robust if and only if it can be written $L(\hat\sigma^2,h)=\tilde C(h)+B(\hat\sigma^2)+C(h)(\hat\sigma^2-h)$ with $C$ strictly decreasing and $\tilde C$ its antiderivative (Proposition~2, printed p.~13).

MSE, $(\hat\sigma^2-h)^2$, is the only robust loss that depends on the forecast error alone, and QLIKE, $\log h+\hat\sigma^2/h$, the only one that depends on the standardised error alone (Proposition~3, printed p.~14). The homogeneous robust losses form a one-parameter family that contains MSE ($b=0$) and QLIKE ($b=-2$) (Proposition~4, printed p.~14) \citep[UTS QFRC Research Paper~175, April 2006, secs.~2--3, printed pp.~4--14]{wide-patton-robust-loss}.

\qffield{Notation} $h_t$ the forecast of the conditional variance $\sigma^2_t$; $\hat\sigma^2_t$ the proxy with $E_{t-1}[\hat\sigma^2_t]=\sigma^2_t$.

\qffield{Conditions} The result needs a conditionally unbiased proxy; it concerns rankings, not the level of expected loss, which noise raises for every forecast (printed p.~4). ``Robust'' here does not mean robust to outliers (n.~6).

\qffield{Practice} With a squared-return proxy, $\hat\sigma^2=\sigma^2\chi^2_1$, MSE and QLIKE both rank the true variance first among forecasts scaled by 0.455, 0.8, 1 and 1.25, while MAE prefers $0.455\sigma^2$, the proxy's conditional median, so MAE (and MAE on logs or standard deviations, by the same argument) favours a forecast biased downward. This argues for reporting MSE and QLIKE when ranking GARCH, GJR and HAR forecasts against realised variance.

\qffield{Pitfalls} Locators refer to the 2006 working paper (\emph{Journal of Econometrics}, 2011). Hansen and Lunde's 2005 comparison, cited in \hyperref[mod:garch-volatility-forecast]{GARCH(1,1) and EWMA}, used seven loss functions, some of them non-robust in Patton's sense.

\qffield{See also} \hyperref[mod:garch-volatility-forecast]{GARCH(1,1) and EWMA}; \hyperref[meth:multiple-testing-control]{multiple testing} (comparing many forecasts).

\qffield{Further reading} \nolinkurl{wide-hansen-lunde-garch11}; \nolinkurl{s7-realized}.

\end{qfmethod}

\begin{qfmethod}[Testing for several structural breaks at unknown dates: least-squares break estimation, trimmed sup-F tests and the sequential choice of the number of breaks (Bai and Perron)]\label{meth:structural-break-tests}
\qffield{Statement} In a linear regression whose coefficients (all or some) shift at $m$ unknown dates, Bai and Perron estimate the break dates by least squares, the break fractions converging at rate $T$. To test no break against $k$ breaks they use $\sup F_T(k;q)$, the largest $F$ statistic over all break partitions in which breaks are at least a fraction $\epsilon$ apart and from the sample ends (``we need to restrict each break date to be asymptotically distinct and bounded from the boundaries of the sample''), generalising Andrews's single-break test; its limit depends on the regressors and on serial correlation and heterogeneity in the errors (sec.~4, printed pp.~16--17).

To choose the number of breaks they propose a sequential procedure: test $\ell$ against $\ell+1$ breaks with $F_T(\ell+1\,|\,\ell)$ in each segment and add a break wherever the test rejects, stopping at the first non-rejection; information criteria are an alternative, with BIC consistent and AIC tending ``to overestimate the dimension'' (sec.~5, printed pp.~26--27) \citep[MIT Department of Economics Working Paper 95-17, September 1995, abstract and secs.~4--5]{wide-bai-perron}.

\qffield{Notation} $T$ observations; $m$ breaks with fractions $\lambda_i=T_i/T$; $q$ the number of coefficients allowed to shift; $\epsilon$ the trimming fraction.

\qffield{Conditions} The critical values depend on $\epsilon$, $q$ and $k$; the asymptotics need breaks bounded away from each other and from the ends. Whether the sequential procedure uses breaks estimated jointly or one at a time matters for the null distribution (sec.~5).

\qffield{Practice} To report break dates that change a forecast ranking, fix $\epsilon$ and the maximum number of breaks before looking, report the sequential test's choice alongside BIC, and state whether rankings of forecasts (\hyperref[meth:volatility-forecast-loss]{loss functions}) change across the estimated regimes.

\qffield{Pitfalls} The source is the 1995 working paper behind the 1998 \emph{Econometrica} article, not the 2003 \emph{Journal of Applied Econometrics} paper with the dynamic-programming algorithm; formulas are not transcribed.

\qffield{See also} \hyperref[meth:volatility-forecast-loss]{volatility forecast comparison}; \hyperref[mod:garch-volatility-forecast]{GARCH and HAR}.

\qffield{Further reading} \nolinkurl{wide-bai-perron} (the 2003 computation paper); \nolinkurl{wide-hamilton-1989}.

\end{qfmethod}

\begin{qfmethod}[Estimating VaR and ES: historical simulation, filtered historical simulation, parametric dynamic models and Monte Carlo as overlapping choices, and the coverage backtest]\label{meth:var-es-estimation-methods}
\qffield{Statement} Christoffersen's textbook \citep[2nd ed., 2012]{wide-christoffersen-elements} treats the methods as layers, not exclusive categories:
\begin{itemize}
\item \emph{Historical simulation} (HS) applies today's holdings to past prices. It assumes ``the distribution of tomorrow's portfolio returns \dots\ is well approximated by the empirical distribution of the past $m$ observations'' (chap.~2, sec.~2.1, PDF pp.~34--35). His conclusion is that ``HS is problematic for computing VaR'', which motivates dynamic models (chap.~2, overview). Weighted HS, which weights recent observations more, is an extension.
\item \emph{Filtered historical simulation} (FHS) ``combines model-based methods of dynamic variance, such as GARCH, with model-free methods of distribution''. It scales the empirical distribution of standardised returns by the model's current volatility (chap.~6, PDF p.~134).
\item \emph{Monte Carlo} and FHS are both used to simulate a dynamic risk model forward, for VaR and ES at horizons longer than a day, where closed forms are lacking. Monte Carlo ``relies on artificial random numbers'' from the model's assumed distribution, FHS ``uses historical random shocks'' (chap.~8, overview, PDF p.~179).
\end{itemize}
\emph{Backtesting.} With $I_{t+1}$ the indicator of a VaR violation at coverage rate $p$, $T_0$ and $T_1$ the numbers of non-violations and violations, the unconditional coverage test is
\[ LR_{uc}=-2\ln\big[(1-p)^{T_0}p^{T_1}\big/\big((1-T_1/T)^{T_0}(T_1/T)^{T_1}\big)\big]\sim\chi^2_1, \]
rejected at 10 percent above 2.7055. An independence test and a combined conditional-coverage test follow (chap.~13, PDF pp.~303--304).

\qffield{Notation} Losses $L$ (positive numbers are losses). The \emph{confidence level} $c$ (for example 99 percent) and the \emph{exceedance} or \emph{coverage} probability $p=1-c$ (1 percent) name the same tail from opposite ends; Christoffersen's $p$ is the latter. $\mathrm{VaR}_c=\inf\{\ell:P(L\le\ell)\ge c\}$, the lower $c$-quantile of the loss. Expected shortfall is the average of VaR over the tail confidence levels,
\[ \mathrm{ES}_c=\frac1{p}\int_c^1\mathrm{VaR}_u\,du=\frac1{p}\Big(E\big[L\,\mathbf 1\{L>\mathrm{VaR}_c\}\big]+\mathrm{VaR}_c\big(P(L\le\mathrm{VaR}_c)-c\big)\Big), \]
the loss-side form of Acerbi and Tasche's tail mean and expected shortfall \citep[Definition~2.6, p.~5; Proposition~3.2, p.~6; eq.~(3.3), p.~7, stated for a profit $X$ and tail level $\alpha=p$]{acerbi-tasche-2001}. The second term gives the atom at VaR exactly the weight needed for the tail to carry probability $p$. $m$ is the HS window.

\emph{Why not ``the expected loss given that VaR is exceeded''.} That phrase equals ES only when $L$ has no atom at $\mathrm{VaR}_c$, for example a continuous distribution; Acerbi and Tasche show that the tail conditional expectations can then differ from ES and fail sub-additivity (pp.~3--5 and Example~5.4). Example: a loss of 100 with probability 0.005, otherwise 0. At $c=99$ percent ($p=1$ percent), $\mathrm{VaR}=0$ because $P(L\le0)=0.995\ge0.99$. Then $\mathrm{ES}=(0.005\times100+0\times(0.995-0.99))/0.01=50$, while $E[L\mid L>\mathrm{VaR}]=100$ and $E[L\mid L\ge\mathrm{VaR}]=0.5$.

In \emph{historical simulation} the same definition is the average of the worst $np$ losses of $n$, with a fractional weight on the last one when $np$ is not an integer: with $n=250$ at 99 percent, $np=2.5$ and $\mathrm{ES}=(\ell_{(1)}+\ell_{(2)}+0.5\,\ell_{(3)})/2.5$, not the mean of the worst two or three.

\qffield{Conditions} HS needs no distributional assumption but assumes the window's returns are draws from tomorrow's distribution: a stationary, independent sample over $m$ days. FHS needs a correct variance model and i.i.d.\ standardised returns. Monte Carlo needs a fully specified distribution. The $\chi^2$ backtest is asymptotic and assumes independent violations under the null.

\qffield{Practice} Synthetic check. Returns are simulated from a GARCH(1,1) with Student-$t$ shocks over 20{,}000 days, and 1 percent HS and FHS VaR are backtested, both on a 250-day window.
\begin{itemize}
\item HS violations occur on 1.62 percent of days, with $LR_{uc}=66$. Their probability the day after a violation is 5.8 percent: they cluster.
\item FHS violations occur on 1.38 percent of days, with $LR_{uc}=25$, the 250-day empirical 1-percent quantile being imprecise. They do not cluster after a violation.
\end{itemize}
A second check verifies the discrete example above, the equality of the integral and fractional-mass forms, historical simulation with 1{,}000 and 250 observations, and agreement with the normal closed form $\varphi(z_{0.99})/0.01=2.665$ for a continuous law. Whether the methods are useful for dynamically updating risk assessments therefore needs this kind of evidence: forecasts with better coverage and independence. That a method can be updated is not enough. Here the evidence is only synthetic, in a world where the GARCH model is true.

\qffield{Pitfalls} The book's ES backtesting and stress-testing chapters, the regulatory ES standard and Kupiec's 1995 paper (\nolinkurl{wide-kupiec-backtest}) are not covered here; extreme value methods for the far tail are in \hyperref[meth:extreme-value-tail-estimation]{extreme value tail estimation}. Page numbers are PDF pages; the printed numbering is offset by a varying amount.

\qffield{See also} \hyperref[mod:factor-risk-model]{factor risk model}; \hyperref[def:common-investor-contagion]{common-investor contagion} (VaR-based leverage in August 2007).

\qffield{Further reading} \nolinkurl{s7-qrm}; \nolinkurl{wide-kupiec-backtest}.

\end{qfmethod}

\begin{qfmethod}[Extreme value theory for tail risk: block maxima and peaks over threshold, the generalized Pareto distribution, the Hill estimator, threshold choice, VaR and ES, and the two-step filtered method]\label{meth:extreme-value-tail-estimation}
\qffield{Statement} Extreme value theory estimates the far tail of a loss distribution from its largest observations. Normalized block maxima of independent losses can converge only to the generalized extreme value family, $H_\xi(x)=\exp(-(1+\xi x)^{-1/\xi})$: Fr\'echet for $\xi>0$, Gumbel for $\xi=0$, Weibull for $\xi<0$ \citep[chap.~7, sec.~7.1.1, pp.~265--266]{mcneil-frey-embrechts-qrm-2005}. Wilmott outlines both maxima and threshold models \citep[chap.~19, sec.~19.8.1, p.~341]{wide-wilmott-pwoqf}. Block maxima waste data, so practice uses \emph{peaks over threshold}. For a wide class of distributions, the excess over a high threshold $u$ tends to a generalized Pareto distribution (GPD), $G_{\xi,\beta}(y)=1-(1+\xi y/\beta)^{-1/\xi}$ \citep[chap.~12, sec.~12.5.1, p.~271]{risk-management-financial-institutions-2023}. Its shape $\xi$ is the same as the GEV shape for the maxima (QRM, Theorem~7.20, pp.~277--278).

\emph{Estimation and risk measures.} Fit $\xi$ and $\beta$ by maximum likelihood to the $n_u$ excesses over $u$, often near the 95th percentile. Then the tail and the risk measures follow:
\[ P(L>x)=\frac{n_u}{n}\Big(1+\frac{\xi(x-u)}\beta\Big)^{-1/\xi}\ (x>u),\qquad \mathrm{VaR}_q=u+\frac\beta\xi\Big[\Big(\frac{n}{n_u}(1-q)\Big)^{-\xi}-1\Big],\qquad \mathrm{ES}_q=\frac{\mathrm{VaR}_q+\beta-\xi u}{1-\xi} \]
(RMFI, secs.~12.5.2--12.5.6, pp.~272--273, eqs.~(12.7)--(12.10); QRM, eqs.~(7.18)--(7.19), p.~283). The \emph{Hill} estimator assumes a regularly varying tail $\bar F(x)=L(x)x^{-\alpha}$ with $\xi=1/\alpha$. It averages the log excesses of the $k$ largest observations over a threshold order statistic (QRM, sec.~7.2.4, pp.~286--289, eq.~(7.23)).

\emph{Two-step (filtered) method.} Returns are serially dependent. McNeil and Frey fit a GARCH model by quasi-maximum likelihood, apply the GPD to the standardized residuals and scale the residual VaR and ES by tomorrow's forecast volatility: $\mathrm{VaR}^t_q=\mu_{t+1}+\sigma_{t+1}q_q(Z)$ (QRM, sec.~7.2.6, pp.~291--292).

\qffield{Notation} $\xi$ shape (tail index $\alpha=1/\xi$); $\beta$ scale; $u$ threshold; $n_u$ exceedances among $n$ losses; $k$ number of upper order statistics. Here $H_k=\tfrac1k\sum_{i=1}^k\ln X_{(n-i+1)}-\ln X_{(n-k)}$, with the $(k+1)$-th largest observation as threshold. QRM's eq.~(7.23) uses the $k$-th largest, and notes that some versions use the $(k+1)$-th. For $\xi>0$, moments of order $1/\xi$ and above are infinite; most financial data give $\xi$ between 0.1 and 0.4 (RMFI, p.~271).

\qffield{Conditions} The GPD is a limit: at any finite threshold the excesses are only approximately GPD, so estimates carry a threshold bias as well as sampling error. A Student $t_4$ has $\xi=0.25$ only asymptotically (QRM, Example~7.29, p.~293). At its 95th percentile the population Hill functional is 0.3489, and the GPD closest to the true excesses has $\xi=0.19$. Thresholds trade bias against variance. In a $t_4$ study QRM finds Hill best for $\xi$ at about 20--30 order statistics of 1{,}000. The GPD quantile estimator is the more robust to threshold choice for the 99 percent VaR (sec.~7.2.5, pp.~289--291). Serial dependence makes maximum likelihood a quasi-likelihood with too-small standard errors (pp.~278--279). For the left tail of returns, change the sign (RMFI, p.~273).

\qffield{Practice} \emph{Reproduced}: RMFI's example, 25 of 500 losses above $u=160$ with $\beta=110.46$ and $\xi=0.354$. It gives $P(L>300)=0.0176$, 99 and 99.9 percent VaR of 399.6 and 1{,}094.6, and ES of 702.0 and 1{,}778.1 (pp.~274--275). The printed figures need $\xi=0.3541$ before rounding; with $\xi=0.354$ exactly the 99.9 percent VaR is 1{,}094.3.

\emph{Checked} (synthetic, $k=n/20$):
\begin{itemize}
\item Exact Pareto data with $\xi=0.25$: over 2{,}000 samples of 5{,}000, Hill averages 0.2505 with standard deviation 0.0157, against $\xi/\sqrt k=0.0158$.
\item $t_4$ with $n=10^6$: Hill is 0.3496 (standard error 0.0014), matching the finite-threshold value 0.3489 and not 0.25.
\item For a Pareto tail the formulas above reproduce the exact VaR and ES. The GPD fit's score vanishes at the estimate.
\end{itemize}
\emph{Reported}: on a 5{,}000-draw $t_4$ sample Hill is 0.339 and the GPD shape 0.200. The 99.5 percent VaR and ES estimates are 4.40 and 5.98, against true values of 4.60 and 6.32. Across 200 samples, Hill's 95 percent bootstrap interval covered the finite-threshold value 190 times and 0.25 never. Over 500 exact-GPD samples of 250 excesses the fitted shape averaged 0.237: maximum likelihood is biased low in small samples. On GARCH(1,1) returns with $t_6$ shocks, the two-step 99 percent one-day VaR was 0.0137 against a true 0.0137, and ES 0.0169 against 0.0176.

\qffield{Pitfalls} Testing an estimated $\xi$ against an asymptotic value at a finite threshold. Reading a Hill plot with no stable region as evidence for one tail index. Applying the unconditional method to clustered returns and reading the standard errors at face value. Extrapolating far beyond the data without stating the threshold and the number of exceedances.

\qffield{See also} \hyperref[meth:var-es-estimation-methods]{VaR and ES estimation}; \hyperref[meth:portfolio-var-methods-decomposition]{portfolio VaR}; \hyperref[mod:arch-conditional-variance]{ARCH and GARCH}.

\qffield{Further reading} \nolinkurl{w2-hill-estimator}; \nolinkurl{w2-mcneil-frey}; \nolinkurl{wide-embrechts-extremal}; \nolinkurl{s7-qrm} (2015 edition).

\end{qfmethod}
